\subsection{乘积与直和的张量积}

设 \((\mathbf{E}_{\lambda})_{\lambda \in L}\) 是右 A-模的一个族， \((\mathbf{F}_{\mu})_{\mu \in M}\) 是左 A-模的一个族，并考虑乘积模 \(C = \prod_{\lambda \in L} E_{\lambda}\) , \(D = \prod_{\mu \in M} F_{\mu}\) 。映射 \(((x_{\lambda}), (y_{\mu})) \mapsto (x_{\lambda} \otimes y_{\mu})\) 把 \(C \times D\) 映入乘积 Z-模

\[
\prod_ {(\lambda , \mu) \in L \times M} (E _ {\lambda} \otimes_ {A} F _ {\mu})
\]

中是 Z-双线性的，并且显然满足条件（1）（第 1 号）。因此存在（第 1 号，命题 1）一个 Z-线性映射，称为典范

\[
f: \left(\prod_ {\lambda \in L} E _ {\lambda}\right) \otimes_ {A} \left(\prod_ {\mu \in M} F _ {\mu}\right)\rightarrow \prod_ {(\lambda , \mu) \in L \times M} (E _ {\lambda} \otimes_ {A} F _ {\mu})\tag{22}
\]

，使得 \(f((x_{\lambda}) \otimes (y_{\mu})) = (x_{\lambda} \otimes y_{\mu})\) 。

当 \(C = R^{L}\) , \(D = S^{M}\) （其中 R（相应地 S）是右（相应地左）A-模）时，典范映射（22）把每个张量积 \(u \otimes v\) （其中 u 是从 L 到 R 的映射，v 是从 M 到 S 的映射）对应到映射 \((\lambda, \mu) \mapsto u(\lambda) \otimes v(\mu)\) （从 \(L \times M\) 到 \(R \otimes_{A} S\) ）；即使在这种情况下，典范映射（22）一般也既不是单射也不是满射（习题3；参见命题7的推论3）。

当 \(E_{\lambda}\) 是 \(((B_{i}^{\prime}); A, (C_{j}^{\prime}))\) -多重模而 \(F_{\mu}(A, (B_{h}^{\prime\prime}); (C_{k}^{\prime\prime}))\) -多重模时，同态（22）也是关于两边的 \(((B_{i}^{\prime}), (B_{h}^{\prime\prime}); (C_{j}^{\prime}), (C_{k}^{\prime\prime}))\) -多重模结构的一个同态。

现在考虑子模 \(\mathbf{E} = \bigoplus_{\lambda \in L} \mathbf{E}_{\lambda}\) （相应地 \(\bigoplus_{\mu \in M} F_{\mu}\) ）（属于 C（相应地 D））；典范单射 \(\mathbf{E} \to \mathbf{C}\) , \(\mathbf{F} \to \mathbf{D}\) 典范地定义了一个 Z-线性映射 \(\mathbf{E} \otimes_{\mathbf{A}} \mathbf{F} \to \mathbf{C} \otimes_{\mathbf{A}} \mathbf{D}\) ，它与映射 (22) 复合后给出一个 Z-线性映射 \(g\) （从 \(\mathbf{E} \otimes_{\mathbf{A}} \mathbf{F}\) 到 \(\prod_{\lambda, \mu} (\mathbf{E}_{\lambda} \otimes_{\mathbf{A}} F_{\mu})\) ），使得

\[
g ((x _ {\lambda}) \otimes (y _ {\mu})) = (x _ {\lambda} \otimes y _ {\mu});
\]

此外，由于族 \((x_{\lambda})\) 和 \((y_{\mu})\) 具有有限支撑，因此 \((x_{\lambda} \otimes y_{\mu})\) 亦然，从而最终 g 是一个典范同态

\[
g: \left(\bigoplus_ {\lambda \in L} E _ {\lambda}\right) \otimes_ {A} \left(\bigoplus_ {\mu \in M} F _ {\mu}\right)\rightarrow \bigoplus_ {(\lambda , \mu) \in L \times M} (E _ {\lambda} \otimes_ {A} F _ {\mu}),\tag{23}
\]

它在与 (22) 相同的条件下是一个多重模同态。

命题7. 典范映射（23）是双射的。

为证明这一点，只需定义一个 \(\mathbf{Z}\) -线性映射 \(h\) （从直和 \(\mathbf{G} = \bigoplus_{(\lambda, \mu) \in L \times M} (\mathbf{E}_{\lambda} \otimes_{\mathbf{A}} \mathbf{F}_{\mu})\) 到 \(\mathbf{E} \otimes_{\mathbf{A}} \mathbf{F}\) ），使得 \(g \circ h\) 和 \(h \circ g\) 都是恒等映射。但是，要定义一个 \(\mathbf{Z}\) -线性映射（从 \(\mathbf{G}\) 到 \(\mathbf{E} \otimes_{\mathbf{A}} \mathbf{F}\) ），只需（§ 1，第 6 号，命题 6）定义一个 \(\mathbf{Z}\) -线性映射

\[
h _ {\lambda \mu}: \mathrm{E} _ {\lambda} \otimes_ {\mathrm{A}} \mathrm{F} _ {\mu} \rightarrow \mathrm{E} \otimes_ {\mathrm{A}} \mathrm{F}
\]

对于每个有序对 \((\lambda, \mu)\) ，并且我们取 \(h_{\lambda \mu} = i_{\lambda} \otimes j_{\mu}\) ，其中 \(i_{\lambda}: \mathbf{E}_{\lambda} \to \mathbf{E}\) 和 \(j_{\mu}: \mathbf{F}_{\mu} \to \mathbf{F}\) 是典范单射。那么显然 \(h \circ g\) 与恒等映射在形如 \(\left(\sum_{\lambda} x_{\lambda}\right) \otimes \left(\sum_{\mu} y_{\mu}\right)\) 的元素上一致，这些元素生成 \(\mathbf{Z}\) -模 \(\mathbf{E} \otimes_{\mathbf{A}} \mathbf{F}\) ；类似地 \(g \circ h\) 与恒等映射在形如 \(\sum_{\lambda, \mu} (x_{\lambda} \otimes y_{\mu})\) 的元素上一致，这些元素生成 \(\mathbf{Z}\) -模 \(\mathbf{G}\) ，因为对每个有序对 \((\lambda, \mu)\) ，乘积

\[
x _ {\lambda} \otimes y _ {\mu} \quad (x _ {\lambda} \in \mathrm{E} _ {\lambda}, y _ {\mu} \in \mathrm{F} _ {\mu})
\]

生成 Z-模 \(E_{\lambda} \otimes_{A} F_{\mu}\) 。由此命题得证。

设 \(u_{\lambda}: \mathbf{E}_{\lambda} \to \mathbf{E}_{\lambda}', v_{\mu}: \mathbf{F}_{\mu} \to \mathbf{F}_{\mu}'\) 是 A-同态；显然图表

\[
\begin{array}{c} \left(\bigoplus_ {\lambda} \mathrm{E} _ {\lambda}\right) \otimes_ {\mathrm{A}} \left(\bigoplus_ {\mu} \mathrm{F} _ {\mu}\right) \longrightarrow \bigoplus_ {\lambda , \mu} (\mathrm{E} _ {\lambda} \otimes_ {\mathrm{A}} \mathrm{F} _ {\mu}) \\ (\oplus u _ {\lambda}) \otimes (\oplus v _ {\mu}) \Bigg \downarrow \qquad \qquad \qquad \qquad \qquad \qquad \qquad \qquad \qquad \qquad \qquad \qquad \qquad \qquad \qquad \qquad \qquad \qquad \qquad \oplus (u _ {\lambda} \otimes v _ {\mu}) \\ \left(\bigoplus_ {\lambda} \mathrm{E} _ {\lambda} ^ {\prime}\right) \otimes_ {\mathrm{A}} \left(\bigoplus_ {\mu} \mathrm{F} _ {\mu} ^ {\prime}\right) \longrightarrow \bigoplus_ {\lambda , \mu} (\mathrm{E} _ {\lambda} ^ {\prime} \otimes_ {\mathrm{A}} \mathrm{F} _ {\mu} ^ {\prime}) \end{array}
\]

是交换的。

推论1. 若左 A-模 F 有一个基 \((b_{\mu})_{\mu \in \mathbb{M}}\) ，则 \(\mathbf{E} \otimes_{\mathbf{A}} \mathbf{F}\) 的每个元素都可以唯一地写成形式 \(\sum_{\mu} (x_{\mu} \otimes b_{\mu})\) ，其中 \(x_{\mu} \in \mathbf{E}\) 且族 \((x_{\mu})\) 具有有限支撑。 \(\mathbf{Z}\) -模 \(\mathbf{E} \otimes_{\mathbf{A}} \mathbf{F}\) 同构于 \(\mathbf{E}^{(\mathbb{M})}\) （视为 \(\mathbf{Z}\) -模）。

基 \((b_{\mu})\) 定义了 \(\mathbf{F}\) 到 \(\bigoplus_{\mu \in \mathbb{M}}\mathrm{Ab}_{\mu}\) 的一个同构，由此（根据命题7）存在一个同构 \(\mathbf{E} \otimes_{\mathbf{A}}\mathbf{F} \to \bigoplus_{\mu \in \mathbb{M}}(\mathbf{E} \otimes_{\mathbf{A}}\mathbf{Ab}_{\mu})\) ；由于 \(\xi \mapsto \xi b_{\mu}\) 是从 \(\mathbf{A}_s\) 到 \(\mathbf{Ab}_{\mu}\) 的同构， \(x \mapsto x \otimes b_{\mu}\) 是从 \(\mathbf{E}\) 到 \(\mathbf{E} \otimes_{\mathbf{A}} (\mathbf{Ab}_{\mu})\) 的同构（根据第 4 号命题4），由此得出推论。如果 \(\mathbf{E}\) 是一个 \(((\mathbf{B}_i^{\prime}); \mathbf{A}, (\mathbf{C}_j^{\prime}))\) -多重模，则典范同构 \(\mathbf{E} \otimes_{\mathbf{A}} \mathbf{F} \to \mathbf{E}^{(\mathbf{M})}\) 是一个 \(((\mathbf{B}_i^{\prime}); (\mathbf{C}_j^{\prime}))\) -多重模同构。

特别地，若 E 也具有基 \((a_{\lambda})_{\lambda \in L}\) ，则每个 \(z \in E \otimes_{A} F\) 可以以唯一的方式写成形式 \(\sum_{\lambda, \mu} (a_{\lambda} \xi_{\lambda \mu}) \otimes b_{\mu}\) ，其中 \(\xi_{\lambda \mu}\) 属于 A（并构成一个有限支撑族）；映射 \(z \mapsto (\xi_{\lambda \mu})_{(\lambda, \mu) \in L \times M}\) 是从 \(E \otimes_{A} F\) 到 \(A^{(L \times M)}\) 的同构（对于 Z-模结构，甚至对于 A 的中心上的模结构）。更特殊地：

推论 2. 若 E 和 F 是交换环 C 上的两个自由模，且 \((a_{\lambda})\) （相应地 \((b_{\mu})\) ）是 C-模 E（相应地，F）的一组基，则 \((a_{\lambda} \otimes b_{\mu})\) 是 C-模 E \(\otimes_{C} F\) 的一组基。

出于语言习惯的滥用，基 \((a_{\lambda} \otimes b_{\mu})\) 有时称为基 \((a_{\lambda})\) 和 \((b_{\mu})\) 的张量积。

注记 (1). 设 E 是自由右 A-模，F 是自由左 A-模， \((a_{\lambda})_{\lambda \in L}\) 是 E 的一组基， \((b_{\mu})_{\mu \in M}\) 是 F 的一组基。每个元素 \(z \in E \otimes_{A} F\) 都可以唯一地写成 \(\sum_{\lambda} a_{\lambda} \otimes y_{\lambda}\) ，其中 \(y_{\lambda} \in F\) ，并且也可以唯一地写成 \(\sum_{\mu} x_{\mu} \otimes b_{\mu}\) ，其中 \(x_{\mu} \in E\) 。如果我们写 \(y_{\lambda} = \sum_{\mu} \eta_{\lambda \mu} b_{\mu}\) , \(x_{\mu} = \sum_{\lambda} a_{\lambda} \xi_{\lambda \mu}\) ，其中 \(\xi_{\lambda \mu}\) 和 \(\eta_{\lambda \mu}\) 属于 A，则 \(\xi_{\lambda \mu} = \eta_{\lambda \mu}\) （对所有 \((\lambda, \mu)\) ），因为

\[
\sum_ {\lambda} \left(a _ {\lambda} \otimes \left(\sum_ {\mu} \eta_ {\lambda \mu} b _ {\mu}\right)\right) = \sum_ {\lambda , \mu} \left((a _ {\lambda} \eta_ {\lambda \mu}) \otimes b _ {\mu}\right) = \sum_ {\mu} \left(\left(\sum_ {\lambda} a _ {\lambda} \eta_ {\lambda \mu}\right) \otimes b _ {\mu}\right).
\]

推论3. 设 \((\mathbf{E}_{\lambda})_{\lambda \in L}\) 是一族右 A-模，F 是一个自由（相应地，有限生成自由）左 A-模。则典范映射 (22)

\[
\left(\prod_ {\lambda \in L} \mathbf {E} _ {\lambda}\right) \otimes_ {\mathbf {A}} \mathbf {F} \rightarrow \prod_ {\lambda \in L} \left(\mathbf {E} _ {\lambda} \otimes_ {\mathbf {A}} \mathbf {F}\right)
\]

是单射（相应地，双射）。

若 \((b_{\mu})\) 是 \(F\) 的一个基，则 \(\left(\prod_{\lambda \in L} E_{\lambda}\right) \otimes_{A} F\) 的每个元素都可以唯一地写成 \(z = \sum_{\mu} ((x_{\lambda}^{(\mu)}) \otimes b_{\mu})\) （推论1）；说它的典范像是零，就意味着对所有 \(\lambda \in L\) , \(\sum (x_{\lambda}^{(\mu)} \otimes b_{\mu}) = 0\) ，因此 \(x_{\lambda}^{(\mu)} = 0\) （对所有 \(\lambda \in L\) 以及所有 \(\mu\) ）（推论1），从而 \(z = 0\) 。

证明当 F 具有有限基时典范映射是双射的，这件事根据命题7立即归结为以下情形：

\(F = A_{s}\) ；但此时两边都被典范地等同于 \(\prod_{\lambda \in L} E_{\lambda}\) （第 4 号，命题4），且在这些等同之后，典范映射 (22) 变为恒等映射。

推论4. 设 A 是一个无零因子的环，E 是一个自由右 A-模，F 是一个自由左 A-模。则关系 \(x \otimes y = 0\) （在 \(\mathbf{E} \otimes_{\mathbf{A}} \mathbf{F}\) 中）蕴含 \(x = 0\) 或 \(y = 0\) 。

设 \((a_{\lambda})\) 是 E 的一组基， \((b_{\mu})\) 是 F 的一组基，并令 \(x = \sum_{\lambda} a_{\lambda} \xi_{\lambda}, y = \sum_{\mu} \eta_{\mu} b_{\mu}\) ；于是 \(x \otimes y = \sum_{\lambda, \mu} ((a_{\lambda} \xi_{\lambda} \eta_{\mu}) \otimes b_{\mu})\) ，且关系 \(x \otimes y = 0\) 蕴含 \(\xi_{\lambda} \eta_{\mu} = 0\) （对每一对有序指标 \((\lambda, \mu)\) ）（推论1）。因此，若 \(x \neq 0\) ，即 \(\xi_{\lambda} \neq 0\) （对至少一个 \(\lambda\) ），则可得 \(\eta_{\mu} = 0\) （对所有 \(\mu\) ），因此 \(y = 0\) 。

推论5. 设 E 为右 A-模，F 为左 A-模，M 为 E 的子模，N 为 F 的子模。若 M 是 E 的直因子且 N 是 F 的直因子，则典范同态 \(M \otimes_{A} N \to E \otimes_{A} F\) 是单射，且 \(M \otimes_{A} N\) 在该同态下的像是 Z-模 \(E \otimes_{A} F\) 的一个直因子。

这直接由命题7得出。

注意，如果 E 是一个 \(((B_{i}^{\prime}); A, (C_{j}^{\prime}))\) -多重模，F 为 \((A, (B_{h}^{\prime\prime}); (C_{k}^{\prime\prime}))\) -多重模，且 M 和 N 是这些多重模中的直因子，则 \(M \otimes N\) 是 \(((B_{i}^{\prime}), (B_{h}^{\prime\prime}); (C_{j}^{\prime}), (C_{k}^{\prime\prime}))\) -多重模 \(E \otimes F\) 的一个直因子。

推论6. 设 P 为投射左 A-模，E、F 为两个右 A-模。对每个单同态 \(u: \mathbf{E} \to \mathbf{F}\) ，同态

\[
u \otimes \mathrm{l} _ {\mathbf {P}}: \mathrm{E} \otimes_ {\mathbf {A}} \mathrm{P} \rightarrow \mathrm{F} \otimes_ {\mathbf {A}} \mathrm{P}
\]

是单射。

存在一个左 A-模 \(\mathbf{Q}\) ，使得 \(\mathbf{L} = \mathbf{P} \oplus \mathbf{Q}\) 是自由的（§2，第 2 号，命题4），且 \(u \otimes 1_{\mathbf{L}}\) （根据命题7）等同于

\[
(u \otimes 1 _ {\mathbf {P}}) \oplus (u \otimes 1 _ {\mathbf {Q}});
\]

因此只需在 P 为自由模的情形下证明该推论（§ 1，第 6 号，命题 7 的推论 1）。同样的论证把问题归结为 \(P = A_{s}\) 的情形，这由第 4 号命题4直接得出。

推论7. 设 C 为交换环。若 E 和 F 是两个投射 C-模，则 E ⊗\_C F 是投射 C-模。

这由推论5以及两个自由C-模的张量积是自由C-模（推论2）这一事实直接得出。

注记2. 在命题7的假设下，设 \(\mathbf{E}_{\lambda}^{\prime}\) 是 \(\mathbf{E}_{\lambda}, \mathbf{F}_{\mu}^{\prime}\) 中前者的子模（后者则是 \(\mathbf{F}_{\mu}\) 的子模），并设 \(\mathbf{E}^{\prime} = \bigoplus_{\lambda \in \mathbf{L}} \mathbf{E}_{\lambda}^{\prime}, \mathbf{F}^{\prime} = \bigoplus_{\mu \in \mathbf{M}} \mathbf{F}_{\mu}^{\prime}\) 。设 \(\operatorname{Im}(\mathbf{E}^{\prime} \otimes_{\mathbf{A}} \mathbf{F}^{\prime})\)

（相应地 \(\operatorname{Im}(\mathbf{E}_{\lambda}^{\prime}\otimes_{\mathbf{A}}\mathbf{F}_{\mu}^{\prime})\) ）表示 \(\mathbf{E}^{\prime}\otimes_{\mathbf{A}}\mathbf{F}^{\prime}\) （相应地 \(\mathbf{E}_{\lambda}^{\prime}\otimes_{\mathbf{A}}\mathbf{F}_{\mu}^{\prime}\) ）在典范映射下在 \(\mathbf{E}\otimes_{\mathbf{A}}\mathbf{F}\) （相应地 \(\mathbf{E}_{\lambda}\otimes_{\mathbf{A}}\mathbf{F}_{\mu}\) ）中的像；则同构 (23) 把子 \(\mathbf{Z}\) -模

\[
\operatorname{Im} \left(\mathrm{E} ^ {\prime} \otimes_ {\mathbf {A}} \mathrm{F} ^ {\prime}\right) \quad \text { and } \quad \bigoplus_ {(\lambda , \mu) \in \mathbf {L} \times \mathbf {M}} \operatorname{Im} \left(\mathrm{E} _ {\lambda} ^ {\prime} \otimes_ {\mathbf {A}} \mathrm{F} _ {\mu} ^ {\prime}\right);
\]

这直接由图表的交换性得出。

\[
\begin{array}{c} \left(\bigoplus_ {\lambda} \mathrm{E} _ {\lambda} ^ {\prime}\right) \otimes_ {\mathrm{A}} \left(\bigoplus_ {\mu} \mathrm{F} _ {\mu} ^ {\prime}\right) \longrightarrow \left(\bigoplus_ {\lambda} \mathrm{E} _ {\lambda}\right) \otimes_ {\mathrm{A}} \left(\bigoplus_ {\mu} \mathrm{F} _ {\mu}\right) \\ \Biggl \downarrow \\ \bigoplus_ {\lambda , \mu} (\mathrm{E} _ {\lambda} ^ {\prime} \otimes_ {\mathrm{A}} \mathrm{F} _ {\mu} ^ {\prime}) \longrightarrow \bigoplus_ {\lambda , \mu} (\mathrm{E} _ {\lambda} \otimes_ {\mathrm{A}} \mathrm{F} _ {\mu}) \end{array}
\]

其中竖直箭头是典范同构。

